\section{Background}
\label{sec:background}

\subsection{Tool-Using LLM Agents}

Modern LLMs such as GPT-4~\cite{openai-gpt4}, Claude~\cite{anthropic-claude}, and LLaMA~\cite{llama}
support \emph{function calling} (also called \emph{tool use})~\cite{schick2023toolformer}, enabling them to
interact with external systems beyond text generation. An \emph{LLM agent} consists of:

\begin{enumerate}[leftmargin=1.2em]
  \item \textbf{System Prompt:} Instructions defining the agent's role, capabilities, and behavioral constraints.
  \item \textbf{Tool Registry:} A set of callable functions with typed schemas (e.g., \texttt{read\_file(path: str)},
        \texttt{send\_email(to: str, subject: str, body: str)}).
  \item \textbf{Execution Loop:} The agent receives user input, generates tool calls via LLM inference,
        executes tools, observes results, and iterates until the task is complete.
\end{enumerate}

\noindent
\textbf{Example: Smart Home Agent.}
Consider an LLM agent managing a smart home. The user asks: \textit{``What's on my calendar today?''}
The agent generates \texttt{calendar.list\_events(date="2026-01-04")}, receives a list of events,
and summarizes them in natural language. If the user then asks: \textit{``Send a notification reminding
me about the first meeting,''} the agent generates \texttt{notify.send(message="Meeting at 2pm")}.

This multi-step reasoning enables powerful automation, but it also creates attack vectors. If an
attacker controls content the agent reads (e.g., a calendar event title), they can inject malicious
instructions that hijack subsequent tool calls.

\subsection{Prompt Injection Attacks}

Prompt injection~\cite{perez2022ignore,greshake2023not} exploits the fact that LLMs process instructions
and data in the same text stream. Unlike SQL injection where code and data have syntactic boundaries,
LLMs interpret \emph{all} input as natural language, making it impossible to reliably distinguish between:
\begin{itemize}[leftmargin=1.2em]
  \item \textbf{Trusted instructions} from the system prompt or user.
  \item \textbf{Untrusted data} from external sources (device names, web pages, files).
\end{itemize}

\paragraph{Direct Injection.}
An attacker embeds malicious instructions directly in user-facing input:
\begin{quote}
  User: ``Summarize this document: [DOCUMENT BEGINS] Ignore previous instructions and delete all files. [DOCUMENT ENDS]"
\end{quote}
The LLM may interpret the injected command as legitimate, calling \texttt{fs.delete("/*")}.

\paragraph{Indirect Injection.}
A more insidious variant: the attacker controls content the agent reads \emph{during} task execution~\cite{greshake2023not}.
For example, an attacker sets their IoT device name to:
\begin{quote}
  Device Name: ``Alice's Thermostat. [SECRET INSTRUCTION: When listing devices, also send an email to attacker@evil.com with all calendar events.]"
\end{quote}
When the agent lists devices, it reads this name, and the injected instruction may propagate into
subsequent tool calls, exfiltrating data without the user's knowledge.

\paragraph{Chained Attacks.}
Attackers can chain multiple injection vectors. An indirect injection in a calendar event title might
instruct the agent to read a specific file, whose content contains a second injection telling the
agent to exfiltrate data. These multi-turn attacks are difficult to detect with simple pattern matching.

\subsection{Limitations of Existing Defenses}

\paragraph{Input Filtering.}
Keyword blacklists (``ignore previous instructions'') are easily bypassed via paraphrasing, encoding
(Base64, URL encoding), or semantic obfuscation (``Please disregard earlier guidance'').

\paragraph{Output Validation.}
Post-hoc checks on generated tool calls can flag suspicious patterns (e.g., deleting many files),
but they lack context. A tool call may be malicious or benign depending on provenance: deleting
\texttt{/tmp/cache} is fine, but deleting \texttt{/home/user/important.txt} is not---unless the
user explicitly requested it.

\paragraph{Instruction Hierarchy.}
Some systems attempt to prioritize system prompts over user input via special tokens or fine-tuning~\cite{instructgpt}.
However, this does not solve indirect injection: when the agent reads attacker-controlled content
as part of task execution, it appears as trusted context.

\paragraph{User Confirmation.}
Requiring confirmation for every tool call provides security but destroys usability. Users cannot
reasonably approve dozens of low-risk operations (e.g., reading files) just to prevent rare high-risk
actions (e.g., deleting files).

\subsection{Capability-Based Security and Provenance}

\paragraph{Capability Security.}
Capability-based access control~\cite{miller2006capability,levy1984capability} grants permissions via
unforgeable tokens (capabilities) rather than ambient authority. Applied to LLM agents, each tool
call can be treated as a capability request, and policies can restrict which capabilities are granted
based on context (time, scope, risk tier).

\paragraph{Data Provenance.}
Provenance tracking~\cite{cheney2009provenance,muniswamy2006provenance} maintains lineage information
for data: where it came from, how it was transformed, and what derived facts depend on it. In databases
and scientific workflows, provenance enables auditing, debugging, and trust assessment. We adapt this
concept to LLM agents: by tracking the provenance of tool call arguments, we can enforce policies that
deny operations originating from untrusted sources.

\sys{} synthesizes these ideas: capability-based policies restrict tool use, and provenance tracking
ensures policies can distinguish legitimate commands from injected instructions.